\section{MỞ ĐẦU}
\subsection{Cơ sở hình thành đề tài}
    Hiểu rõ các hiện tượng tự nhiên và kỹ thuật: Lực điện và điện trường đóng vai trò quan trọng trong nhiều hiện tượng tự nhiên (sét, hiện tượng nhiễm điện do cọ xát) cũng như các ứng dụng công nghệ hiện đại (máy lọc bụi tĩnh điện, công nghệ sơn tĩnh điện, máy in laser, các linh kiện vi điện tử).\\
    
    Giải quyết bài toán thực tế nhiều vật: Trong kỹ thuật và khoa học vật liệu, việc tính toán lực tương tác giữa nhiều hạt mang điện là bước cơ bản để mô phỏng cấu trúc nguyên tử, phân tử hoặc các hệ thống vi cơ điện tử (MEMS).\\
    
    Tuy nhiên, khi số lượng điện tích tăng lên (n ≥ 3), việc tính toán bằng tay trở nên phức tạp, tốn thời gian và dễ dẫn đến sai sót do các phép biến đổi đại số vectơ. Do đó, việc ứng dụng các công cụ tính toán ký hiệu (Symbolic Math) trên các ngôn ngữ lập trình như MATLAB hoặc Python giúp mô phỏng, tự động hóa và trực quan hóa các thành phần lực Fx, Fy một cách nhanh chóng và chính xác.\\
\subsection{Mục tiêu của đề tài}
    \begin{itemize}[leftmargin=1.2cm]
    \item Xây dựng thuật toán và chương trình (bằng MATLAB hoặc Python) ứng dụng phép toán hình thức (symbolic) để tính toán lực tương tác tĩnh điện.
    \end{itemize}
    \begin{itemize}
    \item Xử lý bài toán hệ n điện tích điểm trong mặt phẳng Oxy: Nhập tọa độ ($x_i, y_i$)
    và điện tích qiá trị, tự động xuất biểu thức và giá trị các thành phần lực $F_x, F_y$ tác dụng lên từng điện tích.
    \item Trực quan hóa hệ điện tích và vectơ lực tổng hợp tác dụng lên từng điện tích bằng đồ họa.
    \end{itemize}
\subsection{Ý nghĩa thực tiễn của đề tài}
    \begin{itemize}[leftmargin=1.2cm]
    \item Cung cấp công cụ học tập và kiểm tra kết quả nhanh chóng cho sinh viên khi làm bài tập Điện học.
    \end{itemize}
    \begin{itemize}[leftmargin=1.2cm]
    \item Rèn luyện kỹ năng lập trình mô phỏng vật lý, tư duy thuật toán và sử dụng thư viện đại số ký hiệu (SymPy / Symbolic Toolbox).
    \end{itemize}
\subsection{Phạm vi giới hạn của đề tài}
    \begin{itemize}[leftmargin=1.2cm]
    \item Hệ điện tích được xét là các điện tích điểm đặt cố định trong môi trường chân không (môi trường đẳng hướng, hằng số điện môi ε = 1).
    \end{itemize}
    \begin{itemize}[leftmargin=1.2cm]
    \item Không không gian khảo sát hạn chế trong mặt phẳng tọa độ hai chiều Oxy.
    \end{itemize}